\chapter{Panels, policy changes and banking economics}\label{ch:panels}
\section{Repeated measurement is not a design}
Bank panels combine households, firms, accounts, branches and calendar dates.
They separate persistent differences from changes over time but do not
automatically remove time-varying confounding. In
\begin{equation}
 Y_{it}=\alpha_i+\lambda_t+\beta D_{it}+u_{it},
 \label{eq:twfe}
\end{equation}
unit effects absorb time-constant additive heterogeneity and time effects
absorb common additive shocks. Neither absorbs a local recession affecting
both rollout and outcomes. The coefficient must match the economic target.
Modern designs include group-time effects, imputation, synthetic controls
and interactive factors \citep{panelreview2024,roth2023}.

\subsection{What the within transformation actually removes}
For a balanced panel define
$\widetilde Z_{it}=Z_{it}-\bar Z_{i\cdot}-\bar Z_{\cdot t}
+\bar Z_{\cdot\cdot}$. Applying this linear operation to
Equation~\eqref{eq:twfe} removes both additive effects. A scalar coefficient
is estimated by
\begin{equation}
 \widehat\beta=
 \frac{\sum_{it}\widetilde D_{it}\widetilde Y_{it}}
      {\sum_{it}\widetilde D_{it}^{2}}.
 \label{eq:within}
\end{equation}
The denominator shows the source of information: treatment variation left
after unit and date effects. A national policy adopted by every bank at the
same time has no such variation in this specification. Millions of accounts
cannot rescue that rank failure. For an unbalanced panel, successive simple
demeaning does not generally equal the exact joint projection; a suitable
least-squares or converged absorption routine is needed.

A sufficient model-based condition for the usual within moment is strict
exogeneity, $\E[u_{it}\mid D_{i1},\ldots,D_{iT},\alpha_i]=0$,
together with the appropriate common-effect specification. This excludes
feedback from an unobserved outcome shock into later treatment. Bank
managers often change offers after unexpected withdrawals, so the condition
can fail even if today's treatment predates today's recorded outcome.
Predeterminedness is weaker and does not in general justify the same within
estimator. The policy-study route should therefore start with its design,
not infer exogeneity from a familiar regression formula.

The transformation also cannot remove an arbitrary function of latent
heterogeneity. If untreated deposits are
$Y_{it}(0)=\exp(\alpha_i+f_t)+\epsilon_{it}$, subtracting unit and
date means leaves interactions between the customer and the macro state.
A machine learner fitted to demeaned data inherits that problem. The panel
DML discussion in \citet{paneldml2024} is relevant precisely because
nonlinear latent heterogeneity and dependence complicate otherwise useful
orthogonal-learning constructions.

\section{Staggered adoption and explicit aggregation}
Let $G_i$ be the first treatment date, with $G_i=\infty$ for never-treated
entities. An absorbing rollout defines
\[
 \ATT(g,t)=\E[Y_{it}(g)-Y_{it}(\infty)\mid G_i=g],\qquad t\geq g.
\]
With no anticipation and conditional parallel trends between cohort $g$ and
valid controls $C$,
\begin{align}
 \ATT(g,t)=\E\big[&
 \E[Y_t-Y_{g-1}\mid G=g,X]\nonumber\\
 &-\E[Y_t-Y_{g-1}\mid C=1,X]\mid G=g\big].
 \label{eq:gtatt}
\end{align}
Controls remain untreated over the comparison interval and have support for
the treated cohort. The outer expectation targets the treated covariate
distribution. Never-treated and not-yet-treated comparison groups imply
different assumptions and available dates \citep{callaway2021}.

Aggregation is part of the estimand:
$\theta=\sum_{g,t}w_{gt}\ATT(g,t)$ for declared nonnegative weights
summing to one. Equal customer, branch and treated-period weights answer
different questions. Event-time effects also change cohort composition when
long follow-up is available only for early adopters.

\citet{bacon2021} decompose TWFE into two-by-two comparisons.
\citet{sun2021} show that lead/lag coefficients can mix effects across
event times. \citet{borusyak2024} make target weights and the untreated
outcome model explicit through imputation. The product should expose these
alternatives rather than label every fixed-effects regression valid DiD.

\begin{example}[Perfect untreated trends, misleading TWFE]
Three equal cohorts are observed over four dates. The early cohort adopts
at date 2 with effects $(1,4,8)$; the late cohort adopts at date 4 with
effect $1$; the third never adopts. Untreated outcomes equal a cohort
intercept plus a common linear trend. The four-treated-cell average is
$(1+4+8+1)/4=3.5$.

Residualizing treatment on cohort and time effects yields underlying effect
weights $(0.3,0.3,-0.1,0.5)$. Thus TWFE is
$0.3+1.2-0.8+0.5=1.2$. Comparisons with the never-treated cohort recover
every group-time effect and their average $3.5$.
\end{example}

\begin{figure}[htbp]\centering
\includegraphics[width=\textwidth]{figures/panel_bias.pdf}
\caption{Exact synthetic panel calculation. Contamination occurs despite correct untreated parallel trends.}
\label{fig:panel}
\end{figure}

This is a target and comparison problem; clustering standard errors does
not repair it. A rollout with gradually growing adoption can exhibit this
pattern.

\subsection{A doubly robust cohort comparison}
For one cohort and comparison interval, retain only cohort $g$ and controls
that are untreated over the interval. Let $D=1$ indicate cohort $g$,
$\Delta Y=Y_t-Y_{g-1}$, $p(X)=\Prob(D=1\mid X)$, and
$m_0(X)=\E[\Delta Y\mid D=0,X]$ in this restricted population.
Write $q=\Prob(D=1)$. A population score for the cohort ATT is
\begin{equation}
 \frac{1}{q}\E\left[
 D\{\Delta Y-m_0(X)\}
 -(1-D)\frac{p(X)}{1-p(X)}\{\Delta Y-m_0(X)\}
 \right].
 \label{eq:did-dr}
\end{equation}
Under no anticipation, overlap and conditional parallel trends, the first
term compares treated changes with their counterfactual untreated change.
The second corrects a fitted untreated-change regression using reweighted
control residuals. Its weights map control covariates to the treated
distribution. With either the relevant propensity or untreated-change
regression correct, the mean has the desired robustness; a sample version
also estimates the cohort share and carries its uncertainty. This is a
simple form of the identification and estimation logic developed by
\citet{callaway2021}.

Notice what differs from a static treatment-effect model. The identifying
assumption concerns untreated \emph{changes}, not the levels of untreated
outcomes. A treated branch can have a permanently different customer mix
and deposit level. It becomes a problem when its untreated change would
differ after the rollout, conditional on the stated covariates. Including
post-rollout customer composition as a control can itself remove a policy
mechanism, since the rollout may attract a different customer population.

The common score representation is useful for software. Once each
$\ATT(g,t)$ has cluster contributions, an aggregate effect uses the
same declared weights on those contributions. A covariance matrix over
event times then supports simultaneous bands and financial aggregation.
Repeatedly calling a package for separate dates and discarding cross-date
covariance loses information required to value a multi-month policy.

\subsection{The comparison set is an economic choice}
Never-treated branches may be structurally unlike rollout branches.
Not-yet-treated branches can be more comparable but can anticipate a
scheduled change. A one-month anticipation allowance changes the last valid
baseline and excludes controls close to their own adoption. When every
branch eventually adopts, late effects may have no untreated controls.
The honest output is then a shorter identified horizon or a more model-
dependent extrapolation, rather than a graph with silently declining support.

Event-time aggregation needs a second decision. At event month 12 only early
cohorts may remain observable. Comparing that average with event month 1
can conflate effect dynamics with a different cohort mix. A balanced event
window estimates the same cohorts at every horizon but discards information.
The report should show both the target weights and the set of contributing
cohorts. A bank sponsor can then decide whether the retained target
population matches the planned rollout.

\section{Trend sensitivity and reversals}
An insignificant pretrend test does not establish parallel trends. Low power
and post-selection affect interpretation. \citet{rambachan2023} allow
specified smoothness or relative-magnitude deviations. Show conclusions
over a justified sensitivity range and the point at which the decision
changes.

\subsection{Sensitivity gives the assumption an economic scale}
Suppose a simple two-period estimate is $\widehat\tau$. Let
$\delta$ be the treated-minus-control change in untreated potential
outcomes, so that $\E[\widehat\tau]=\tau+\delta$ under the rest of
the design. If substantive reasoning limits $\delta$ to $[-M,M]$,
the compatible treatment effects shift by that range. An approximate
sampling interval $[L,U]$ becomes a conservative union
$[L-M,U+M]$ when that bound is treated as fixed. This elementary
construction is deliberately separate from the more structured event-study
confidence sets in \citet{rambachan2023}.

The parameter should be expressed in business units. If a positive policy
decision reverses when untreated deposits could have grown only USD\,20
per eligible customer more rapidly in treated branches, that is useful
information for a funding economist. A sensitivity parameter whose
relation to deposit growth is hidden behind a default slider is less
useful. Historical untreated variation, known local shocks and institutional
rollout rules can motivate a range, but do not prove its bound.

Pre-treatment coefficients have two roles: they can reveal contradictions
with a design, and they can help calibrate a sensitivity analysis. Passing
a pretrend test is not a license to assume $M=0$. Selection of the final
design after searching for quiet pretrends can distort the subsequent
interval. The product records the original study plan and exploratory
alternatives so that the final inference can be interpreted honestly.

Bank treatments also change intensity or reverse: rates move repeatedly,
contact policies stop and restart, and branch services expand in stages.
Absorbing binary-adoption methods may be inappropriate.
\citet{dechaisemartin2024}, in the archived May 2026 revision, study
intertemporal effects with status-quo comparisons for groups sharing initial
treatment. No anticipation, comparison support and treatment-path meaning
remain essential. Preserve intensity and version rather than collapsing
every history to first treatment.

\section{Dynamic panels and latent factors}
In
\[
 Y_{it}=\rho Y_{i,t-1}+\beta A_{it}+\alpha_i+\varepsilon_{it},
\]
demeaning correlates the transformed lagged outcome and disturbance in a
short panel. Differencing removes $\alpha_i$ but not endogeneity. Under
suitable serial-correlation and initial-condition restrictions, lagged
levels instrument differenced equations. Moments such as
$\E[Y_{i,t-s}\Delta\varepsilon_{it}]=0$, $s\geq2$, require justification.
\citet{arellano1991} develop relevant specification tests. First-order
correlation in differenced errors may be mechanical; higher-order tests,
instrument validity and sensitivity to instrument count matter.

\subsection{Deriving an internal-instrument moment}
First differencing the dynamic model gives
\[
 \Delta Y_{it}=\rho\Delta Y_{i,t-1}+\beta\Delta A_{it}
              +\Delta\varepsilon_{it}.
\]
The regressor $\Delta Y_{i,t-1}$ contains $\varepsilon_{i,t-1}$,
whereas the disturbance contains $-\varepsilon_{i,t-1}$.
Even independent level errors therefore create a correlation. Under serially
uncorrelated level innovations and suitable initial conditions, $Y_{i,t-2}$
is orthogonal to both $\varepsilon_{it}$ and $\varepsilon_{i,t-1}$.
It can instrument $\Delta Y_{i,t-1}$. Older levels provide additional
moments, but the allowable lags for $A$ depend on whether actions are
strictly exogenous, predetermined or contemporaneously endogenous.

Stack admissible instruments in $Z_i$ and differenced residuals in
$v_i(\vartheta)$. The sample moment is
$\bar g(\vartheta)=N^{-1}\sum_i Z_i'v_i(\vartheta)$, and GMM
minimizes $\bar g(\vartheta)'W\bar g(\vartheta)$ for a suitable
weighting matrix. Robust covariance estimation uses unit-level moment
contributions; the two-step weighting matrix and its finite-sample
behavior need explicit treatment. This construction explains why dynamic
panel software requires the instrument set, transformation and inference
convention to be part of the saved model \citep{arellano1991}.

Persistence creates another problem. If levels are highly persistent,
lagged levels may be weak instruments for changes. Enlarging the instrument
set can overfit endogenous changes and make overidentification diagnostics
unhelpful. A sensible comparison limits lag depth, considers collapsed
instrument sets and reports sensitivity; it does not select the specification
with the most reassuring diagnostic $p$-value. A system of additional
level-equation moments needs additional stationarity or initial-condition
restrictions. It is a different identifying model, not an automatic
upgrade whenever difference GMM is imprecise.

Interactive-factor models
$Y_{it}(0)=\alpha_i+\lambda_t+\ell_i^\top f_t+\varepsilon_{it}$
permit different exposures to common shocks. Factor rank, untreated-data
availability, treatment selection and extrapolation govern their use.
An estimated factor is not automatically a measured confounder.

Panel DML combines flexible nuisances with panel structure, but nonlinear
latent heterogeneity need not disappear under a within transformation.
\citet{paneldml2024} examine point estimation in simulation.
\citet{cluster2021} establish multiway-cluster DML results under specified
exchangeability and dissociation conditions. These are not guarantees
for arbitrary serial dependence.

\section{Synthetic controls and structural time series}
For a small number of policy changes, construct
$\widehat Y_{1t}(0)=\sum_{j\in\mathcal D}w_jY_{jt}$ from suitable
untreated donors, often with nonnegative weights summing to one.
Pre-fit assesses one aspect of plausibility; interpolation bias and
spillovers remain concerns \citep{abadie2010,abadie2021}.
Outcome-model augmentation can correct imperfect fit while implying
negative weights \citep{benmichael2021}.

Synthetic DiD combines unit and time weighting \citep{synthdid2021}.
Its inspected package documents block treatment: paper extensions do not
prove implementation support. Bayesian structural time-series methods offer
another counterfactual forecast \citep{brodersen2015}. Controls must be
unaffected and their counterfactual relationship must remain credible.

Nearby branches receiving displaced customers after a closure may be
invalid donors. National campaign controls may share exposure. Placebo
dates, donor sensitivity and pre-period holdouts can challenge a design;
they cannot create an untreated comparison where none exists.

\subsection{What a donor fit identifies}
One useful model writes untreated outcomes as
$Y_{it}(0)=\alpha_i+\lambda_t+\ell_i'f_t+\epsilon_{it}$.
If donor weights reproduce the treated unit's intercept and factor loading,
their weighted untreated mean reproduces its systematic counterfactual.
A long and informative pre-period can constrain those loadings through
outcome fit. It does not ensure that a novel post-period factor lies in
the same span or that donor outcomes remain unaffected.

The conventional nonnegative simplex constraint controls extrapolation:
the synthetic unit lies in the donors' convex hull. If the treated unit is
outside that hull, poor fit can persist even with many donors. Augmentation
uses an outcome model to correct residual imbalance and can permit negative
effective weights \citep{benmichael2021}. It exchanges a stricter interpolation
restriction for a model-dependent extrapolation. Synthetic DiD balances
both units and dates in its main block-adoption setting
\citep{synthdid2021}. These are distinct responses to imperfect comparison
data; neither is uniformly preferable.

Inference must match the design. A placebo distribution across donors can
be informative when units are suitably exchangeable under the null and
fit differences are handled. It is not automatically a randomized test
when a regulator deliberately selected one troubled institution.
Time-series posterior predictive intervals condition on the counterfactual
model and its unaffected controls \citep{brodersen2015}. Their coverage is
not a test of the causal exclusion. A bank report should distinguish
sampling, model and design uncertainty instead of merging all three into
one shaded band labeled ``confidence.''

\section{Financial interpretation and uncertainty}
Independent policy changes constrain information more than raw account
counts. \citet{bertrand2004} demonstrate the consequences of ignored
serial correlation. Preserve assignment clusters and disclose few-cluster
or short-time limitations.

The lending design of \citet{mian2006} compares bank supply among
multi-bank borrowers; it does not directly estimate all borrowers' losses.
The mortgage equilibrium analysis of \citet{benetton2021} adds demand
and supply assumptions. Regional effects in \citet{beraja2019} cannot
simply be summed into an aggregate monetary-policy effect. Preserve the
distinction between an identified local effect, a transported effect and
a structural scenario.
